\subsubsection{Non-Clifford Gates}\label{sec:non-clifford-gates}

While the Hadamard and CNOT gates are examples of Clifford gates, there are many other gates that do not have this property. The most important example of a non-Clifford gate is the $T$ gate.

\highspace
The \definitionWithSpecificIndex{T Gate}{Single Qubit Gates: T Gate}{} (also called the $\pi/8$ gate) is a \textbf{single-qubit gate that applies a specific phase shift to the state} \ket{1}, while leaving \ket{0} \textbf{unchanged}. Its matrix representation is:
\begin{equation}
    T =
    \begin{bmatrix}
        1 & 0 \\
        0 & e^{i\pi/4}
    \end{bmatrix}
    =
    \begin{bmatrix}
        1 & 0 \\[.3em]
        0 & \frac{1 + i}{\sqrt{2}}
    \end{bmatrix}
\end{equation}
This means:
\begin{equation}
    T \ket{0} = \ket{0}
    \qquad
    T \ket{1} = e^{i\pi/4} \ket{1}
\end{equation}
On the Bloch sphere, the $T$ gate corresponds to a rotation around the $Z$-axis by an angle of $\pi/4$ (or 45 degrees). In terms of the Pauli operators, the $T$ gate transforms them as follows:
\begin{equation}
    \begin{aligned}
        T X T^\dagger &=
        \frac{1}{\sqrt{2}} (X + Y) \\
        T Y T^\dagger &=
        \frac{1}{\sqrt{2}} (Y - X) \\
        T Z T^\dagger &=
        Z
    \end{aligned}
\end{equation}
Since the $T$ gate transforms the $X$ and $Y$ operators into linear combinations of $X$ and $Y$ (which are not single Pauli operators), the $T$ gate is \textbf{not a Clifford gate.}

\highspace
\begin{deepeningbox}[: The $T$ Gate is the Square Root of the $S$ Gate]
    The $T$ gate is the square root of the Phase gate $S$ (which is a Clifford gate) because applying $T$ twice gives exactly $S$.

    \highspace
    If we apply the $T$ gate twice, we have:
    \begin{equation*}
        T^2 =
        \begin{bmatrix}
            1 & 0 \\
            0 & e^{i\pi/2}
        \end{bmatrix}
    \end{equation*}
    Where:
    \begin{equation*}
        e^{i\pi/4} \cdot e^{i\pi/4} = \exp\left(
            \frac{i\pi}{4} + \frac{i\pi}{4}
        \right) = 
        \exp\left(
            \frac{2 i\pi}{4}
        \right) =
        \exp\left(
            \frac{i\pi}{2}
        \right) =
        e^{i\pi/2}
    \end{equation*}
    Now using Euler's formula, we can express $e^{i\pi/2}$ as:
    \begin{equation*}
        \begin{aligned}
            e^{i\theta} &= \cos\left(\theta\right) + i \sin\left(\theta\right) \\
            e^{i\pi/2} &= \cos\left(\frac{\pi}{2}\right) + i \sin\left(\frac{\pi}{2}\right) \\
            &= 0 + i \cdot 1 \\
            &= i
        \end{aligned}
    \end{equation*}
    Therefore, we have:
    \begin{equation*}
        T^2 =
        \begin{bmatrix}
            1 & 0 \\
            0 & i
        \end{bmatrix} =
        S
    \end{equation*}
    Finally, we can also express $T$ as the square root of $S$:
    \begin{equation*}
        T =
        \sqrt{S} =
        \begin{bmatrix}
            1 & 0 \\
            0 & \sqrt{i}
        \end{bmatrix}
    \end{equation*}
    In summary, \hl{applying the $T$ gate twice gives us the $S$ gate}, confirming that $T$ is indeed the square root of $S$.
\end{deepeningbox}

\highspace
\begin{flushleft}
    \textcolor{Green3}{\faIcon{question-circle} \textbf{Why are Non-Clifford Gates important?}}\phantomsection\label{question:why-non-clifford-gates-important}
\end{flushleft}
The real question is: \textbf{why can some quantum circuits be simulated efficiently on a classical computer, while others cannot?} To answer this, we must compare two ways of describing a quantum state.
\begin{itemize}
    \item The \textbf{full state vector representation}, where an $n$-qubit quantum state has the form (\autopageref{eq:general-multiple-qubit-state}):
    \begin{equation*}
        \ket{\psi} = \sum_{k=0}^{2^n-1} c_k \ket{k}
    \end{equation*}
    This requires storing $2^n$ complex amplitudes (e.g., 1 qubit requires 2 amplitudes, 2 qubits require 4 amplitudes, etc.), which \textbf{grows exponentially} with the number of qubits.

    \item The \textbf{stabilizer representation}. Certain special states can be described much more compactly. Instead of storing all amplitudes, we store a small set of Pauli operators (stabilizers) that uniquely characterize the state. Only $n$ independent stabilizers are needed for $n$ qubits. For example, 10 qubits can be described by just 10 stabilizers, which is a huge reduction compared to the $2^{10} = 1024$ amplitudes needed in the full state vector representation. This representation \textbf{grows only linearly} with the number of qubits, making it much more efficient for certain states.
\end{itemize}
\hl{Clifford gates use the stabilizer representation}, because transform Pauli operators into other Pauli operators. Therefore, we \hl{start with $n$ stabilizers} ($n$ Pauli operators) that describe the initial state, and \hl{after applying Clifford gates, we still have $n$ stabilizers} that describe the new state. This allows to \hl{leave the state description compact and efficient.}

\highspace
\textcolor{Green3}{\faIcon{question-circle} \textbf{What happens at a Non-Clifford Gate?}} Consider the $T$ gate. It transforms, for example, the $X$ operator into a linear combination of $X$ and $Y$:
\begin{equation*}
    T X T^\dagger =
    \frac{1}{\sqrt{2}} (X + Y)
\end{equation*}
The result is not a single Pauli operator, but a sum of two Pauli operators. After another Non-Clifford transformation, each term may split again. So \textbf{repeated Non-Clifford gates can cause the number of terms to double repeatedly, leading to exponential growth}.

\highspace
\textcolor{Green3}{\faIcon{tools} \textbf{Computational Consequence.}} Classical simulation cost depends on how much information must be stored.
\begin{itemize}
    \item Clifford circuits keep the description small.
    \item Non-Clifford circuits cause the description to grow rapidly.
\end{itemize}
Therefore, if a \hl{circuit contains only Clifford gates, a classical computer can efficiently simulate it.} In contrast, if the circuit includes Non-Clifford gates such as $T$, simulation generally becomes exponentially more difficult.

\highspace
\textcolor{Green3}{\faIcon[regular]{star} \textbf{So, why are Non-Clifford gates important?}} If Clifford circuits are classically simulable, then they cannot provide full quantum computational power. To achieve universal quantum computation, we need gates that escape the stabilizer framework. The $T$ gate does exactly this. By adding the $T$ gate (or any other suitable non-Clifford gate) to the Clifford gate set, we obtain a universal set of quantum gates. This allows us to implement arbitrary quantum algorithms, including computations believed to be infeasible for classical computers.